\chapter{Metodologia}
\section{Enquadramento Epistemológico Metodológico}
A Fundamentação Epistemológiva desta pesquisa baseia - se no paradigma de Design Science Research
(DSR), voltado ao desenvolvimento, avaliação e validação de artefados tecnológicos para resolver 
problemas práticos e gerar novos conhecimentos cientificos. Operacionalmente, a investigação
adota o método Design Science Research Methodology (DSRM) proposto por \citeonline{peffers2007design}
que organiza a pesquisa nas etapas de identificação do problema, definição dos objetivos, projeto
e desenvolvimento, demonstração, avaliação e comunicação.
\section{Unidade Experimental de Fatores Controlados}
A uniddade Esperimental deste trabalho é constituida pelo nó sensor físico baseado no microcontrolador
Heltec Wifi Lora 31 v2, integrado aos sensores de gases (MQ - 9, MQ - 131, MQ 135) e variaveis
microclimáticos (DHT11).

Para avaliar o desempenho de artefato embarcadom dois fatores principais serão manipulados e
testados de forma controlada:

\begin{enumerate}
    \item \textbf{Concorrencia e Escalonamento do Firmware:} Avaliados de dois niveis de comparaçãp:
                                                             nó operando com RTOS (escalonamento preemptivo 
                                                             por prioridade de tarefas) versus nó sem RTOS(execução
                                                             sequencial não deterministico)
    \item \textbf{Niveis de Quantização do Modelo TinyML:}  Submetido a 5 estratégias de otimização pós - treinamento
                                                            com TensorFlow Lite Micro, baseando - se em \citeonline{rodrigues2024tinyml}
                                                            Float32 (padrão) Float16, Faixa Dinâmica (Dynamic Range), Full Int8 (inteiro
                                                            completo) e Weight Only (apenas peso)
\end{enumerate}

\section{Variaveis de Resposta}

Os efeitos das variações dos fatores controlados serão medidos por duas dimensões de variáveis de Resposta

\subsection{Desempenho de Computação, Hardware e Rede}
\begin{itemize}
    \item \textbf{Latência Computacional ($\mu$s):} Tempo de processamento consumido durante a inferência
                                                    do modelo na borda;
    \item  \textbf{Packet Delivery Ratio (PDR\%)}: Taxa percentual de pacotes entregues com sucesso vis LoRaWan;
    \item  \textbf{Throughput de Transmissão e Recepção (b/s)}: Volumetria de bits transmitidos e recebidos;
\end{itemize}

\subsection{Acurácia da Inteligencia Artificial}

\begin{itemize}
    \item \textbf{Man Square Error (MSE)};
    \item \textbf{Root Mean Square Error (RMSE)};
    \item \textbf{Mean Absolut Error ($R^2$)};
\end{itemize}

\section{Procedimentos Metodológicos}

O clico de engenharia e avaliação do artefato organiza - se em cinco etapas sequenciais

\begin{enumerate}
    \item \textbf{Etapa 1: Análise de Requisitos e Problematização}: Revisão sistemática da literatura
                                                                     (2020 - 2026), especificação de requisitos e modelagem da 
                                                                     arquitetura distribuida;
    \item \textbf{Etapa 2: Projeto do Sistema Embarcado}: Desenvolvimento do hardware do nó sensor e codificação
                                                          em C/C++ do firmware concorrente gerenciado pelo FreeRTOS
    \item \textbf{Etapa 3: Desenvolvimento do Pipeline TinyML}: Treinamento offline de redes neurais
                                                                em python, quantização pós - treinamento e conversão em array C/C++
                                                                para execução via TensorFlow Lite Micro;
    \item \textbf{Etapa 4: Integração do Middleware}: Configuração do LoRaWan com servidor ChirpStack e publicação contextual via plataforma reescrita
                                                      com persistencia em PostgreSQL;
    \item \textbf{Etapa 5: Validação em Campo}: Implantação piloto contínua por mais de 30 dias 
                                                em ambientes internos e externos do campus amazônico para coleta in loco
                                                de métricas computacionais e ambientais;
\end{enumerate}



\section{Protocolo de Controle e Coleta de Dados}
\subsection{Controle Experimental e Calibração Metrológica}
\begin{itemize}
    \item \textbf{Isolamento de Condições de Contorno}: Nos teste comparativos entre firmwares 
                                                        (FreeRTOS vs Não-RTOS), os nós sensores serão
                                                        instaldos rigorosamente no mesmo lo
\end{itemize}